\section{Introduction to LoRa in Vehicular Environments}

Most Intelligent Transportation Systems (ITS) research and deployments are centered around Dedicated Short-Range
Communications (DSRC, IEEE 802.11p) and Cellular V2X (C-V2X). These are designed specifically for vehicular
communication and they are still the main focus in standards-oriented deployments \cite{secure_iov,lorawan_v2roadside}.
In this report, we treat LoRa mainly as a low-cost way to experiment with decentralized V2V safety-message
dissemination when the channel is very constrained.

Prior literature suggests LoRa can still work as an experimental platform for low-rate vehicular dissemination
over sub-GHz links, even with mobility, as long as the protocol is designed with the constraints in mind
\cite{exp_study_lora_vehicle,lora_arch_v2x_move,urban_vanets_2023}. LoRa's Chirp Spread Spectrum and high
receiver sensitivity can help with longer-range packet exchange, but the low data rate, duty-cycle limits, and
long Time-on-Air make it important to keep forwarding efficient \cite{lora_mesh_iot,routing_mesh_2020}.
So, most LoRa-V2V work focuses on selective, low-overhead safety dissemination rather than high-throughput
communication.

\section{The Paradigm of Centralized vs. Decentralized V2V}

One thing that comes up a lot in V2X literature is that many practical architectures depend on infrastructure
like RSUs, gateways, or cellular coordination entities \cite{secure_iov,lorawan_v2roadside}. That is fine for
planned deployments, but it does not directly answer a different question we care about here: how much
cooperative safety dissemination is possible when vehicles rely only on peer-to-peer forwarding and local state.

This matters when infrastructure is sparse, temporarily unavailable, or just outside a project-scale budget.
Studies on LoRa mesh and emergency communication show that infrastructure-free operation is feasible, but
practical evidence of fully decentralized vehicular behavior (with realistic forwarding constraints) is still
limited \cite{lora_mesh_emergency,lora_mesh_iot,routing_mesh_2020}.

\section{Evaluating Existing Network and Link Layer Protocols}

At the link layer, LoRa-focused work repeatedly shows that channel access policy matters a lot when there is
lots of broadcast traffic. If access is uncoordinated, collisions go up, but tightly scheduled approaches can
add synchronization overhead that is hard to maintain in ad-hoc mobility \cite{lora_mesh_iot,routing_mesh_2020}.
This pushes designs toward lightweight access control that can run on resource-constrained nodes without
centralized timing.

At the network layer, prior work highlights two common problems in low-bandwidth vehicular meshes: too much
control overhead from route-centric protocols in fast-changing topologies, and broadcast storms from
unrestricted flooding \cite{routing_mesh_2020,vdtn_routing_eval}. Because of this, the literature generally
leans toward managed dissemination with controls like hop limits, duplicate suppression, and forwarding
deferral. The goal is to reduce redundant rebroadcasts while still keeping the dissemination reach
\cite{lora_reliable_v2v_2025,lora_mesh_emergency}.

These findings are the reason our design direction is: lightweight channel-aware transmission at the link
layer and controlled flooding at the network layer, without relying on route tables.

\section{Research Gap and Proposed Architecture}

Based on the reviewed work, the gap we target is that there is limited end-to-end validation of a fully
decentralized LoRa-based V2V mesh system that can operate without RSUs, fixed gateways, or centralized
scheduling/control, especially at a final-year project scale \cite{lora_reliable_v2v_2025,routing_mesh_2020}.

This report addresses that gap using an end-to-end architecture and dual-mode validation, which are presented
in the methodology and results sections \cite{lora_mesh_emergency,lora_reliable_v2v_2025}.

